%% parts/p07/functional-interpretations.tex --- Phase-4 P7 chapter 7.6
%% New content: N-a-053, N-a-054, N-a-055.
%% Term-calculus rules remain in registry/calculi. Dialectica soundness and
%% relative consistency are proved in the linked proof unit. The separate
%% realizability and HA constructive-property results retain their stated status.

\chapter{Functional Interpretations and Constructive Properties}
\label{sec:p07-functional-interpretations}
\chaptermark{Functional Interpretations}

Functional interpretations attach computational witnesses to derivations.
This chapter develops a fixed Dialectica translation, its quantifier-free
target, and its HA soundness proof. It then treats Kleene number realizability
and the distinct disjunction and numerical-existence properties.  The principal sources are
\cite{Godel1958,Troelstra1973,AvigadFeferman1998,Kleene1945,
Troelstra1998,TroelstraVanDalen1988}.

Heyting arithmetic \(\mathsf{HA}\) has the exact standing fixed in
Definition~\ref{def:p07-heyting-arithmetic}; this chapter does not introduce a
second presentation of the theory.

\section{System T and the Dialectica interpretation}
\label{sec:p07-dialectica}

\begin{hbdefinition}{def:p07-system-t}{Godel's System T@G\"odel's System T}
The finite types form the least set containing the ground type \(0\) and
closed under the right-associative type former
\((\alpha,\beta)\mapsto\alpha\to_{\mathrm{ty}}\beta\).  The typed terms form
the least family containing variables of every finite type, the constants
listed below, and the application and typed-lambda-abstraction constructions
of the explicitly typed lambda-calculus on
\hyperref[sheet:stlc]{\texttt{sheet:stlc}}, now with a ground type \(0\), a
constant \(0:0\), successor \(S:0\to_{\mathrm{ty}}0\), and, for every finite type
\(\alpha\), a recursor
\[
 R_\alpha:
 \alpha\to_{\mathrm{ty}}
 (0\to_{\mathrm{ty}}\alpha\to_{\mathrm{ty}}\alpha)
 \to_{\mathrm{ty}}0\to_{\mathrm{ty}}\alpha.
\]
The zero and successor computation rules for this displayed argument order
live only on \hyperref[sheet:system-t]{\texttt{sheet:system-t}}.  That sheet,
rather than this body entry, fixes the operational surface used below.
This is \emph{System T}: a typed term calculus of primitive recursion at all
finite types.  It differs from \(\mathsf{HA}\), which is a first-order theory
of numbers, and from the quantifier-free theory \(\mathsf T_{\mathrm{qf}}\)
below, whose
formulas assert equality of System-T terms.  The distinction is necessary
because the Dialectica theorem extracts System-T terms while proving their
matrix property in \(\mathsf T_{\mathrm{qf}}\).
\index{System T@System \(T\)}
\index{primitive recursive functional!finite type}
\end{hbdefinition}

\begin{hbdefinition}{def:p07-dialectica-target}{The quantifier-free target theory \(\mathsf T_{\mathrm{qf}}\)}
The terms and finite types are those of \cref{def:p07-system-t}.
An atomic formula is an equality of two terms of the same finite type;
formulas are built with the classical propositional connectives.
The theory has classical propositional reasoning and modus ponens,
reflexivity, and the replacement schema
\[
 s=t\land A[s/x]\ \longrightarrow\ A[t/x],
\]
with matching types and capture-free substitution. Its term equations
are the compatible beta equations of
\hyperref[sheet:stlc]{\texttt{sheet:stlc}} and the recursor equations of
\hyperref[sheet:system-t]{\texttt{sheet:system-t}}. Definitional conversion is closed under typed term contexts; no rule
inferring functional equality from arbitrary pointwise equality is added.
Its additional numerical axioms are
\(Sn\ne0\) and \(Sn=Sm\to n=m\).
Same-type term substitution preserves theorems. Open induction permits
one to infer \(A(t)\) from proofs of \(A(0)\) and
\(A(n)\to A(Sn)\), where \(n\) and \(t\) have ground type and
\(A\) is quantifier free; its other typed parameters remain free.
Set \(\bottom\) to \(0=S0\). Addition and multiplication below are
the System-T terms defined by primitive recursion, not new primitives.

This fixes the intensional, all-types-equality version called \(T\)
in Section~2.5 of the held July 20, 1999 author preprint associated with
\cite{AvigadFeferman1998}, in the lambda presentation allowed in its
Section~2.2. It does not add the higher-type equality-testing constants
of \(I\text{-}T\), or replace equality by the different rule of
\(\mathrm{WE}\text{-}T\). Numerical equality tests used below are
constructed from the recursors. The theory is separate from both HA
and the term-reduction calculus: its open induction and propositional
inferences are assertions about equations, not term-reduction steps.
\index{System T@System \(T\)!quantifier-free target theory}
\end{hbdefinition}

\begin{hbdefinition}{def:p07-dialectica-translation}{A fixed Dialectica translation}
For every arithmetic formula \(\varphi\), define a finite-type formula
\[
 \varphi^{\mathrm D}\defeq
 \exists\vec x\,\forall\vec y\,
 \varphi_{\mathrm D}(\vec x,\vec y),
\]
where the witness tuple \(\vec x\) and challenge tuple \(\vec y\) may be
empty and the matrix \(\varphi_{\mathrm D}\) is quantifier-free in
\(\mathsf T_{\mathrm{qf}}\).  Every component of those tuples has the finite
type determined by the recursive clause in which it occurs; bound variables
are chosen fresh and every displayed substitution is capture avoiding.  If
\(\varphi^{\mathrm D}=\exists\vec x\forall\vec y,\varphi_{\mathrm D}\) and
\(\psi^{\mathrm D}=\exists\vec u\forall\vec v,\psi_{\mathrm D}\), the clauses are
\[
\begin{aligned}
 \chi^{\mathrm D}
   &\defeq \chi &&(\chi\text{ atomic}),\\
 (\varphi\land\psi)^{\mathrm D}
   &\defeq \exists\vec x,\vec u\,\forall\vec y,\vec v\,\\
   &\qquad (\varphi_{\mathrm D}(\vec x,\vec y)\land
        \psi_{\mathrm D}(\vec u,\vec v)),\\
 (\varphi\lor\psi)^{\mathrm D}
   &\defeq \exists z,\vec x,\vec u\,\forall\vec y,\vec v\,\\
   &\qquad ((z=0\land\varphi_{\mathrm D}(\vec x,\vec y))\lor
        (z=1\land\psi_{\mathrm D}(\vec u,\vec v))),\\
 (\varphi\to\psi)^{\mathrm D}
   &\defeq \exists\vec f,\vec g\,\forall\vec x,\vec v\,\\
   &\qquad (\varphi_{\mathrm D}(\vec x,\vec g(\vec x,\vec v))\to
        \psi_{\mathrm D}(\vec f(\vec x),\vec v)),\\
 (\forall z\,\varphi(z))^{\mathrm D}
   &\defeq \exists\vec f\,\forall z,\vec y\,
       \varphi_{\mathrm D}(\vec f(z),\vec y,z),\\
 (\exists z\,\varphi(z))^{\mathrm D}
   &\defeq \exists z,\vec x\,\forall\vec y\,
       \varphi_{\mathrm D}(\vec x,\vec y,z).
\end{aligned}
\]
Negation is implication to \(\bottom\), so its clause is derived rather than
added.  The tag \(z\) is necessary to distinguish the selected disjunct;
the functionals \(\vec f\) transform witnesses and \(\vec g\) transform a
purported counterexample to the conclusion into a challenge to the premise.
This is an inductively defined translation of formulas, not formula
equivalence in the source logic.  Its clauses are verified in the
soundness proof below.
\index{Dialectica interpretation}
\index{functional interpretation!Dialectica}
\end{hbdefinition}

\subsection{The extracted functionals and their verification}
\label{sec:p07-dialectica-proof}

All tuples below abbreviate finite lists of typed terms. No product
type or simultaneous-recursion primitive is added to the arrow-only
System T. Every finite type has a closed default term: use zero at
base type, and a constant function to the default of the result type
at an arrow type. Empty witness and challenge tuples require no term.

\begin{hblemma}{thm:p07-t-simultaneous-recursion}{proved=annex:p7.d.1}{Tuple recursion}
For a finite list of types \(\sigma_1,\ldots,\sigma_r\), initial
terms \(a_i:\sigma_i\), and typed step terms
\(F_i(n,x_1,\ldots,x_r):\sigma_i\), System T defines terms
\(X_i(n):\sigma_i\) such that \(\mathsf T_{\mathrm{qf}}\) proves
\[
 X_i(0)=a_i,\qquad
 X_i(Sn)=F_i(n,X_1(n),\ldots,X_r(n)).
\]
Parameters of all finite types are permitted. The term construction
and the proofs of these equations are effective.
\end{hblemma}

\begin{hblemma}{thm:p07-t-matrix-decision}{proved=annex:p7.d.2}{Deciding the matrices}
For every Dialectica matrix of an arithmetic formula there is a
System-T term \(d_A\) of ground type, with the same parameters,
such that the target proves
\[
 d_A=0\quad\longleftrightarrow\quad A.
\]
There are typed conditional terms
\(\operatorname{If}_\sigma(e,u,v)\) taking value \(u\) when
\(e=0\), and \(v\) otherwise. These constructions test only
numerical equalities. They do not decide equality at arbitrary
functional types.
\end{hblemma}

\begin{hblemma}{thm:p07-t-backward-induction}{proved=annex:p7.d.3}{Backward challenges}
Let \(P(n,\vec y)\) be a quantifier-free target formula, with
\(n\) of ground type and a finite typed tuple \(\vec y\).
If \(\mathsf T_{\mathrm{qf}}\) proves
\[
 P(0,\vec y),\qquad
 P(n,\vec b(n,\vec y))\longrightarrow P(Sn,\vec y),
\]
then it proves \(P(n,\vec y)\). The proof transformation uses only
System-T terms, propositional reasoning, and open induction. In
particular it does not instantiate a pointwise induction hypothesis
at a different unlicensed value of its challenge parameter.
\end{hblemma}

\begin{hblemma}{thm:p07-dialectica-induction}{proved=annex:p7.d.4}{The induction schema}
Every instance of the HA induction schema has a Dialectica realizer
in System T, with its matrix proved in
\(\mathsf T_{\mathrm{qf}}\), uniformly in the formula and its
numerical parameters. This treats the full induction axiom of the
handbook's HA presentation, not merely an inference rule for two
already proved closed premises.
\end{hblemma}

The proof of soundness below checks every logical rule, the selected
arithmetic axioms including predecessor existence, and induction.
The conditional for contraction selects a false challenge when one
is available. With the convention \(d_A=0\leftrightarrow A\), it is
\[
 \operatorname{If}(d_{A(\vec x,\vec y_1)},\vec y_2,\vec y_1),
\]
component by component. In the held July 20, 1999 preprint, the first
line of p.~13 prints the last two arguments in the opposite order,
although p.~12 uses the same zero-for-truth convention. That displayed
choice fails when the first challenge makes the matrix true and the
second makes it false. The proof uses the displayed corrected order;
no claim about the wording of the separate 1998 published version is
made from that preprint.


\begin{hbtheorem}{thm:p07-dialectica-soundness}{proved=annex:p7.d.5}{Dialectica soundness for HA}
Let \(\varphi\) be a closed arithmetic sentence.  If
\(\derives_{\mathsf{HA}}\varphi\), then
one can effectively extract a tuple of closed System-T terms \(\vec t\) such that
the quantifier-free target theory proves
\[
  \derives_{\mathsf T_{\mathrm{qf}}}\varphi_{\mathrm D}(\vec t,\vec y).
\]
The challenge variables \(\vec y\) remain free, equivalently universally
closed at the metalevel.
\end{hbtheorem}

The target statement is Theorem~2.4.1 in the held July 20, 1999 author
preprint associated with \cite{AvigadFeferman1998}. The annex proves it
for the explicitly specified target and the handbook's HA axioms.
It is an effective metatheorem about derivations, not internal HA
reflection or an interpretation of arbitrary analysis.

\begin{hbcorollary}{thm:p07-dialectica-consistency}{proved=annex:p7.d.6}{Relative consistency consequence}
In primitive recursive arithmetic, for the proof predicates of the
fixed source and target presentations,
\[
  \mathsf{Con}(\mathsf T_{\mathrm{qf}})\to\mathsf{Con}(\mathsf{HA}).
\]
\end{hbcorollary}

The proof uses the interpretation's fixed treatment of falsity and the
primitive-recursive transformation of proof codes. It does not prove
\(\mathsf{Con}(\mathsf{HA})\) internally in HA or assert absolute
consistency of either theory.

\begin{hbconvention}{def:p07-dialectica-proof-boundary}{Scope of the soundness proof}
The proof constructs the tuple recursors, numerical matrix tests,
contraction challenges, and full induction-axiom realizer in the stated
arrow-only term calculus. It checks all twelve groups of logical rules
in the source presentation, equality, and every arithmetic axiom of the
handbook's HA, including predecessor existence. Its higher-type equality
is the specified intensional equality; no higher-type decision constants
or unrestricted extensionality principle are added. The relative
consistency consequence is a transformation of finite proof codes.
These results concern HA and its chosen quantifier-free target; extensions
to analysis require their own interpreted axioms.
\end{hbconvention}

\section{Kleene number realizability}
\label{sec:p07-number-realizability}

\begin{hbdefinition}{def:p07-number-realizability}{Number realizability}
Fix an acceptable numbering of the unary partial recursive functions, a
primitive-recursive pairing function, and projections \(\pi_0,\pi_1\).  For
a closed arithmetic formula \(\varphi\), the phrase
``the number \(e\) realizes \(\varphi\)'' is defined by structural recursion using
Table~\ref{tab:p07-number-realizability}.  It names a metalinguistic relation
between a natural number and an arithmetic formula; it is not the forcing or
satisfaction relation.  For formulas with numerical parameters, the cited
extension requires one index to realize all closed numeral instances
uniformly; that parameter clause is not silently identified with any one row
of the closed-formula table.  Pairing, program application, and definedness
are used by the soundness and ECT entries below \cite{Troelstra1998}.
\index{realizability!Kleene number}
\index{number realizability}
\end{hbdefinition}

\begingroup\small
\begin{longtable}{@{}p{.19\textwidth}p{.72\textwidth}@{}}
\caption{Kleene number-realizability clauses.}
\label{tab:p07-number-realizability}\\
\toprule
Formula & A number \(e\) realizes it exactly when \\
\midrule
\endfirsthead
\caption[]{Kleene number-realizability clauses (continued)}\\
\toprule
Formula & A number \(e\) realizes it exactly when \\
\midrule
\endhead
\midrule\multicolumn{2}{r}{\emph{continued on next page}}\\
\endfoot
\endlastfoot
Atomic \(\chi\) & \(\chi\) is true in the standard natural numbers; the value of
\(e\) carries no further information \\
\(\varphi\land\psi\) & \(\pi_0(e)\) realizes \(\varphi\) and
\(\pi_1(e)\) realizes \(\psi\) \\
\(\varphi\lor\psi\) & \(\pi_0(e)=0\) and \(\pi_1(e)\) realizes \(\varphi\), or
\(\pi_0(e)=1\) and \(\pi_1(e)\) realizes \(\psi\) \\
\(\varphi\to\psi\) & For every \(a\) realizing \(\varphi\), the indexed
computation with program number \(e\) on input \(a\) terminates and its output
realizes \(\psi\) \\
\(\forall x\,\varphi(x)\) & For every \(n\in\N\), the indexed computation
with program number \(e\) on input \(n\) terminates and its output realizes the
result of substituting the standard numeral denoting \(n\) for \(x\) in
\(\varphi\) \\
\(\exists x\,\varphi(x)\) & \(\pi_0(e)\) is the numerical witness and
\(\pi_1(e)\) realizes the result of substituting the standard numeral denoting
\(\pi_0(e)\) for \(x\) in \(\varphi\) \\
\(\bottom\) and \(\neg\varphi\) & No number realizes \(\bottom\); negation is
\(\varphi\to\bottom\), so \(e\) realizes \(\neg\varphi\) exactly when no
number realizes \(\varphi\), in which case the implication clause is vacuous \\
\bottomrule
\end{longtable}
\endgroup

\begin{hbtheorem}{thm:p07-number-realizability-soundness}{proved=annex:p7.real.1}{Soundness of number realizability for HA}
There is an effective transformation which, from an
\(\mathsf{HA}\)-derivation of a closed arithmetic sentence \(\varphi\), produces a
natural number realizing \(\varphi\).
\end{hbtheorem}

The annex constructs the programs uniformly in numerical parameters and
checks all logical rules, equality, and the complete arithmetic axiom
list. Effective parameterization supplies the required program indices.

\subsection{Formal realizability and extended Church's thesis}
\label{sec:p07-formal-realizability}

The external relation on numbers becomes an arithmetic interpretation
only after its computation language and its treatment of undefined
terms have been fixed. The following presentation specifies both.
The standard-numeral conclusion at the end is a metatheorem about
closed sentences, not an internal reflection principle for arithmetic.
For the wider account of realizability compare \cite{Troelstra1998};
the results here refer to the calculus and clauses displayed below.

\begin{hbdefinition}{def:p07-partial-arithmetic}{Arithmetic with partial application}
First expand HA definitionally by symbols for primitive-recursive
functions, with their defining equations. Use a fixed primitive-recursive
coding of finite computations, a computation predicate \(T(e,x,s)\),
and an output function \(U(s)\). Here \(T\) is represented by its
zero-valued primitive-recursive characteristic function. Choose the
ordinary program coding with effective parameter substitution; the
finite-code constructions used below are performed in this coding.

Adjoin a binary partial application term \(e\mathbin{\cdot}x\).
Write \(t\downarrow\) for \(t=t\), with equality read strictly:
both terms must denote. Variables denote, and functions and atomic
predicates are strict in their arguments. Quantifier instantiation
at \(t\) and existential introduction with witness \(t\) require
\(t\downarrow\). The other intuitionistic natural-deduction rules
retain their usual discharge and eigenvariable conditions. Equality
has variable reflexivity, symmetry, transitivity and the rule from
\(t=u\) and \(C(t)\) to \(C(u)\), with capture-avoiding
substitution. Axioms are universally closed; their instantiation at
partial terms retains the definedness guards. Equality does not assert
that a partial function is defined on equal arguments.
Primitive-recursive functions are total on defined arguments. The application axiom is
\[
 e\mathbin{\cdot}x=y\quad\longleftrightarrow\quad
 \exists s\bigl(T(e,x,s)\land U(s)=y\bigr).
\]
Together with HA's arithmetic axioms and induction for every formula
of this expanded language, these rules and axioms define
\(\mathsf{HA}^{\ast}\). There is no totality axiom for application.

An expression \(\lambda x.t\) in the proofs denotes the code of a
program evaluating \(t\), not a new object-language binder. Its
code is primitive recursive in the other free numerical parameters,
even when evaluating \(t\) may diverge. Application satisfies
\((\lambda x.t)\mathbin{\cdot}a\simeq t[a/x]\), where
\(u\simeq v\) means
\((u\downarrow\lor v\downarrow)\to u=v\).
Pairing and its projections \(\pi_0,\pi_1\) are the fixed total
primitive-recursive operations.
\end{hbdefinition}

\begin{hblemma}{thm:p07-partial-arithmetic-conservative}{proved=annex:p7.ect.1}{Eliminating partial terms}
\(\mathsf{HA}^{\ast}\) is conservative over the original HA for
sentences of the original arithmetic language. Finite program
composition, parameter substitution, and iteration have their
computation equations provable in \(\mathsf{HA}^{\ast}\).
\end{hblemma}

\begin{hbdefinition}{def:p07-formal-number-realizability}{Formal number realizability}
Define \(t\mathrel{\mathrm r}A\) for every term \(t\), including
partial terms, by recursion on the formula \(A\). Every clause
includes \(t\downarrow\). Beyond this conjunct, the clauses are:
\begin{enumerate}
\item an atom \(P\): \(P\) itself;
\item \(A\land B\):
      \(\pi_0t\mathrel{\mathrm r}A\land
        \pi_1t\mathrel{\mathrm r}B\);
\item \(A\lor B\): the tag \(\pi_0t\) is zero or one,
      and its corresponding payload \(\pi_1t\) realizes the
      selected formula;
\item \(A\to B\):
      \(\forall a(a\mathrel{\mathrm r}A\to
                      (t\mathbin{\cdot}a)\mathrel{\mathrm r}B)\);
\item \(\forall xA\):
      \(\forall x((t\mathbin{\cdot}x)\mathrel{\mathrm r}A)\);
\item \(\exists xA\):
      \(\pi_1t\mathrel{\mathrm r}A[\pi_0t/x]\).
\end{enumerate}
The disjunction clause is written without disjunction or existential
quantification in the interpreting formula. Let \(v(d)\) be zero
exactly for \(d=0,1\), and one otherwise. Its exact expansion is
\[
 v(\pi_0t)=0\ \land\
 (\pi_0t=0\to\pi_1t\mathrel{\mathrm r}A)\ \land\
 (\pi_0t=1\to\pi_1t\mathrel{\mathrm r}B).
\]
The arithmetic decidability of the tag proves its equivalence to the
selected-branch description. Falsity is the false atomic equation.
All substitutions are capture avoiding. In particular
\(t\mathrel{\mathrm r}A\to t\downarrow\) holds even when
\(A\) is an implication with an unrealizable antecedent.
\end{hbdefinition}

\begin{hbdefinition}{def:p07-almost-negative-ect}{Restricted extended Church's thesis}
Let \(\mathcal N\) be the formulas built from atoms and falsity
using conjunction, implication, and universal quantification in the
partial-term language. For \(D(x)\in\mathcal N\) and arbitrary
\(B(x,y)\), the partial-term scheme \(\mathsf{ECT}_0\) consists
of the universal closures of
\[
 \forall x\bigl(D(x)\to\exists yB(x,y)\bigr)
 \ \to\
 \exists e\,\forall x\bigl(D(x)\to
       (e\mathbin{\cdot}x\downarrow\land B(x,e\mathbin{\cdot}x))\bigr).
\]
Parameters are allowed; the variables bound in this display are fresh
where necessary. The restriction concerns the antecedent \(D\),
not the conclusion matrix \(B\).

In total primitive-recursive arithmetic, an \emph{almost-negative}
formula has no disjunction and has existential quantifiers only
immediately before atomic formulas. Each such existential atom is
equivalent in \(\mathsf{HA}^{\ast}\) to the definedness of the
program searching for its witness. Consequently its almost-negative
combinations have equivalents in \(\mathcal N\). Eliminating
partial application yields the familiar graph formulation
\[
 \forall x(D(x)\to\exists yB(x,y))\ \to\
 \exists e\,\forall x\bigl(D(x)\to
          \exists y(\chi(e,x,y)\land B(x,y))\bigr),
\]
where \(\chi\) is the displayed computation graph, not a new
primitive function. Conversely, eliminating a partial-term atom
produces a \(\Sigma^0_1\) formula, which can be written as an
existential atomic primitive-recursive equation. Thus the antecedent
classes correspond up to provable equivalence; they are not literally
the same strings in the two languages.
\index{extended Church thesis@extended Church's thesis}
\index{almost-negative formula}
\end{hbdefinition}

\begin{hblemma}{thm:p07-formal-realizability-lemmas}{proved=annex:p7.ect.2}{Canonical realizers and formal soundness}
Every formula \(t\mathrel{\mathrm r}A\) belongs to \(\mathcal N\).
For each \(D\in\mathcal N\) there is a term \(c_D\), with no
additional free parameters, such that \(\mathsf{HA}^{\ast}\) proves
\[
 (t\mathrel{\mathrm r}D\to D)
 \quad\hbox{and}\quad
 (D\to c_D\mathrel{\mathrm r}D).
\]
A derivation of \(A\) in \(\mathsf{HA}^{\ast}\) effectively yields
a term \(t\) and a derivation of \(t\mathrel{\mathrm r}A\) in
that same theory. Every instance of the restricted
\(\mathsf{ECT}_0\) also has a realizer whose realizing property is
provable in the base theory \(\mathsf{HA}^{\ast}\).
\end{hblemma}

\begin{hbtheorem}{thm:p07-ect-characterization}{proved=annex:p7.ect.3}{Formal number-realizability characterization}
For every formula \(A\), the explicitly specified theory satisfies
\[
 \mathsf{HA}^{\ast}+\mathsf{ECT}_0\ \vdash\
 A\longleftrightarrow\exists a(a\mathrel{\mathrm r}A).
\]
For every closed sentence \(A\) in its language,
\[
 \mathsf{HA}^{\ast}+\mathsf{ECT}_0\vdash A
 \quad\Longleftrightarrow\quad
 \mathsf{HA}^{\ast}\vdash\overline n\mathrel{\mathrm r}A.
\]
Here the right-hand side asserts the existence of such a proof for
some standard natural number \(n\). The numeral assertion is external. It does not assert a
primitive-recursive numeral-extraction bound or formalized
arithmetical reflection inside HA.
\end{hbtheorem}

\begin{hbtheorem}{thm:p07-unrestricted-ect-inconsistent}{proved=annex:p7.ect.4}{Unrestricted ECT}
Allowing arbitrary antecedents in the displayed
\(\mathsf{ECT}_0\) scheme makes \(\mathsf{HA}^{\ast}\)
inconsistent, even when its premise is restricted to unique
existence of the output. The contradiction requires no excluded-middle
axiom for arbitrary arithmetic formulas.
\end{hbtheorem}


\begin{hbdefinition}{def:p07-markov-principle}{Markov's principle for decidable arithmetic predicates}
For an arithmetic predicate \(\varphi(x)\), the relevant instance of
\emph{Markov's principle} is the universal closure of
\[
 \forall x\,(\varphi(x)\lor\neg\varphi(x))
 \to
 \bigl(\neg\neg\exists x\,\varphi(x)\to\exists x\,\varphi(x)\bigr).
\]
The displayed decidability premise is part of the scheme.  This is an
object-language principle over arithmetic, distinct from the classical
metatheoretic step used to analyze realizers.
\index{Markov principle@Markov's principle}
\end{hbdefinition}

\subsection{Modified realizability and total extracted functions}
\label{sec:p07-modified-realizability}

Number realizers may be partial programs that terminate only on
realizing inputs. System-T terms instead denote total functionals on
their finite types. This difference supplies a nonderivability argument
for Markov's principle. The interpretation used here is modified
realizability; it is distinct from both ordinary number realizability
and the proof-carrying relation used for HA proof extraction.

\begin{hbdefinition}{def:p07-modified-realizability}{Modified realizability}
Interpret ground type by \(\mathbb N\) and an arrow type by all
set-theoretic functions between the interpreted types. A tuple below
is a finite metalinguistic list; it adds no product type to System T.
Associate a positive witness tuple type \(\rho(A)\) with each
arithmetic formula: an atom has an empty tuple, conjunction concatenates
the two tuple types, disjunction adds a numerical tag and both tuples,
existential quantification adds a numerical witness, and universal
quantification replaces each component type \(\sigma\) by
\(0\to_{\mathrm{ty}}\sigma\). For implication \(A\to B\),
replace each component \(\tau\) of \(\rho(B)\) by the arrow type
from the successive components of \(\rho(A)\) to \(\tau\).
If the input tuple is empty there are no input arrows; if the output
tuple is empty there are no components at all.

Define the semantic relation \(\operatorname{mr}_A\) recursively:
\begin{enumerate}
\item \(\operatorname{mr}_P()\) means that the numerical atom \(P\)
      is true;
\item \(\operatorname{mr}_{A\land B}(\vec a,\vec b)\) means
      \(\operatorname{mr}_A(\vec a)\land\operatorname{mr}_B(\vec b)\);
\item \(\operatorname{mr}_{A\lor B}(d,\vec a,\vec b)\) means
      \((d=0\land\operatorname{mr}_A(\vec a))\lor
      (d=1\land\operatorname{mr}_B(\vec b))\);
\item \(\operatorname{mr}_{A\to B}(\vec F)\) means that for every
      \(\vec a\) of the input types,
      \(\operatorname{mr}_A(\vec a)\) implies
      \(\operatorname{mr}_B(\vec F(\vec a))\);
\item \(\operatorname{mr}_{\forall nA(n)}(\vec F)\) means
      \(\operatorname{mr}_{A(n)}(\vec F(n))\) for every natural
      number \(n\);
\item \(\operatorname{mr}_{\exists nA(n)}(n,\vec a)\) means
      \(\operatorname{mr}_{A(n)}(\vec a)\).
\end{enumerate}
Object-language numerical parameters remain parameters of this relation.
All quantification over witness tuples in these clauses ranges over
the full interpreted finite types, not just denotations of terms.
In particular negation has an empty positive tuple, and its condition
is that no positive witness realizes its argument.
\end{hbdefinition}

\begin{hblemma}{thm:p07-system-t-numerical-computability}{proved=annex:p7.mr.1}{Numerical computation}
Every closed System-T term of ground type has a finite compatible
beta/recursor reduction to a numeral. Such a numeral can be found
effectively and has the term's unique value in the full finite-type
interpretation. Consequently every closed System-T term of type
\(0\to_{\mathrm{ty}}0\) defines a total computable numerical
function. This assertion does not require a claim that every possible
reduction sequence terminates.
\end{hblemma}

\begin{hbtheorem}{thm:p07-modified-realizability-soundness}{proved=annex:p7.mr.2}{Modified soundness}
An HA proof of \(A(\vec n)\) yields witness terms effectively.
They have the component types \(\rho(A)\), with free parameters
among \(\vec n\). For every numerical tuple their denotations satisfy
\(\operatorname{mr}_{A(\vec n)}\). Closed conclusions yield closed
terms. No constants beyond System T are used.
\end{hbtheorem}

\begin{hblemma}{thm:p07-bounded-modified-realizability}{proved=annex:p7.mr.3}{Bounded truth}
For every fixed bounded arithmetic formula \(B(\vec n)\):
\begin{enumerate}
\item any tuple satisfying \(\operatorname{mr}_B\) implies the
      numerical truth of \(B\);
\item System T defines a canonical tuple \(\vec b(\vec n)\),
      total for all parameters, which satisfies
      \(\operatorname{mr}_{B(\vec n)}\) whenever \(B(\vec n)\)
      is true;
\item System T defines a numerical decision term for its truth.
\end{enumerate}
Thus truth is equivalent to existence of a modified realizer for these
bounded formulas. In particular there is a uniform System-T witness
for \(\forall t(B(\vec n,t)\lor\neg B(\vec n,t))\).
\end{hblemma}


\begin{hblemma}{thm:p07-markov-number-realizability}{proved=annex:p7.real.2}{Markov realizability}
In the classical metatheory, ordinary Kleene number realizability
validates every instance of \cref{def:p07-markov-principle}, including
its displayed decidability premise. A uniform realizer searches the
successive outputs of the supplied decision realizer. The double-negation
realizer guarantees that this search terminates; it is not replaced by
an unrestricted termination assertion.
\end{hblemma}

\begin{hbtheorem}{thm:p07-markov-realizability-boundary}{proved=annex:p7.mr.4}{Nonderivability of Markov's principle}
\(\mathsf{HA}\) does not derive the full scheme in
\cref{def:p07-markov-principle}. Already one instance with a bounded
predicate and a universally closed numerical parameter is not provable.
This is a nonderivability result, separate from validation by ordinary
number realizability.
\end{hbtheorem}

The two assertions use different interpretations. Ordinary number
realizability validates the scheme by a search terminating on realizing
inputs. Modified realizability extracts total System-T functions; the
halting argument shows why such a function cannot realize the chosen
Markov instance. Validation in the first semantics is therefore not
mistaken for a countermodel in the second.  Realizability
topoi are outside this row; see \cite{VanOosten2008} only for further reading.

\begingroup\small
\begin{longtable}{@{}p{.25\textwidth}p{.31\textwidth}p{.35\textwidth}@{}}
\caption{Number-realizability applications and their honest conclusions.}
\label{tab:p07-realizability-applications}\\
\toprule
Input & Supported conclusion & Unsupported inference \\
\midrule
\endfirsthead
\caption[]{Number-realizability applications and their honest conclusions (continued)}\\
\toprule
Input & Supported conclusion & Unsupported inference \\
\midrule
\endhead
\midrule\multicolumn{3}{r}{\emph{continued on next page}}\\
\endfoot
\endlastfoot
HA derivation & Effective number-realizer extraction,
Theorem~\ref{thm:p07-number-realizability-soundness} & Truth of every
classical principle \\
Formal realizability plus almost-negative-domain \(\mathsf{ECT}_0\) & The
characterization in Theorem~\ref{thm:p07-ect-characterization} & Unrestricted
ECT, or the informal Church--Turing thesis as an axiom \\
Ordinary number realizability and Markov's principle & Realizability validates
the displayed MP scheme & A realizability countermodel to MP; a different
named construction is required \\
Realizability of \(\bottom\) & No number realizes \(\bottom\), giving the
corresponding relative consistency reading & An internal HA proof of its own
consistency \\
\bottomrule
\end{longtable}
\endgroup

\section{Disjunction and numerical existence}
\label{sec:p07-constructive-properties}

\subsection{Realizers that carry derivations}
\label{sec:p07-proof-carrying-realizability}

An ordinary number realizer need not supply an HA proof. To extract
proofs, we use a different relation and construct its proof components
explicitly. HA's axioms remain unchanged.

\begin{hbdefinition}{def:p07-proof-carrying-realizability}{Proof-carrying realizers}
Fix effective codes for the finite HA derivations used above. For a
closed arithmetic formula \(A\), write \(r\Vdash^{\mathrm p}A\)
when \(r=\langle p,w\rangle\), \(p\) is an HA proof of \(A\),
and the payload \(w\) satisfies the following recursive clauses:
\begin{enumerate}
\item an atomic formula is true in the standard numerical structure;
\item for \(A\land B\), the two projections of \(w\) are
      proof-carrying realizers of \(A\) and \(B\);
\item for \(A\lor B\), the first projection is zero or one and
      the second is a proof-carrying realizer of the selected disjunct;
\item for \(A\to B\), \(w\) is an index taking every
      proof-carrying \(A\)-realizer to a terminating
      proof-carrying \(B\)-realizer;
\item for \(\forall nA(n)\), \(w\) is an index taking each number
      \(n\) to a terminating proof-carrying realizer of
      \(A(\overline n)\);
\item for \(\exists nA(n)\), the first projection of \(w\) is a
      number \(n\), and the second is a proof-carrying realizer of
      \(A(\overline n)\).
\end{enumerate}
Falsity is the false atomic equation, and negation is implication to
falsity. This is an external relation on numbers and formulas. It is
not the ordinary realizability relation with an unproved inference to
provability appended to it. In particular, every recursively required
subrealizer must carry its own proof, not merely a proof of the whole
outer formula. Free numerical parameters are handled by numeral
substitution.
\end{hbdefinition}

\begin{hbtheorem}{thm:p07-proof-carrying-soundness}{proved=annex:p7.real.4}{Soundness with proof components}
From an HA derivation of \(A(\vec n)\), one can effectively construct
a recursive procedure that, on every numerical tuple \(\vec n\),
terminates with an \(r\) satisfying
\(r\Vdash^{\mathrm p}A(\overline{\vec n})\).
Every proof component is obtained by finite HA proof constructions;
no normalization theorem for arithmetic proof terms, truth-reflection
principle, or rule assembling an infinite family of numeral proofs is
assumed.
\end{hbtheorem}


\begin{hblemma}{thm:p07-lj-disjunction-last-rule}{proved=annex:p7.functional.1}{Last-rule lemma}
If there is a derivation on
\hyperref[sheet:lj]{\texttt{sheet:lj}} of the empty-antecedent sequent
\(\Rightarrow\varphi\lor\psi\) containing neither Cut nor Mix, then that sheet
has a derivation of \(\Rightarrow\varphi\) or a derivation of
\(\Rightarrow\psi\).
\end{hblemma}

\Cref{thm:p07-lj-exact-cut-elimination} supplies the complete Cut and
Mix elimination needed before applying this last-rule lemma. Its scope
is the stated logical sheet, without arithmetic or equality axioms.

\begin{hbtheorem}{thm:p07-ipc-disjunction-property}{proved=annex:p7.real.3}{Propositional disjunction property}
For propositional formulas \(\varphi\) and \(\psi\), if
\(\derives_{\mathsf{LJ}}\varphi\lor\psi\), then
\(\derives_{\mathsf{LJ}}\varphi\) or
\(\derives_{\mathsf{LJ}}\psi\), where derivability is in the propositional
fragment of \hyperref[sheet:lj]{\texttt{sheet:lj}} with empty antecedent.
\end{hbtheorem}

\begin{hbtheorem}{thm:p07-ha-disjunction-property}{proved=annex:p7.real.5}{Disjunction property for HA}
For closed arithmetic sentences \(\varphi\) and \(\psi\), if
\(\derives_{\mathsf{HA}}\varphi\lor\psi\), then
\(\derives_{\mathsf{HA}}\varphi\) or
\(\derives_{\mathsf{HA}}\psi\).
\end{hbtheorem}

\begin{hbtheorem}{thm:p07-ha-numerical-existence}{proved=annex:p7.real.6}{Numerical existence property for HA}
Let \(\varphi(x)\) be an arithmetic formula whose only possible free variable is
the displayed number variable.  If
\(\derives_{\mathsf{HA}}\exists x\,\varphi(x)\), then for some \(n\in\N\),
the result of substituting the standard numeral denoting \(n\) for \(x\) in
\(\varphi\) is derivable in \(\mathsf{HA}\).
\end{hbtheorem}

Propositional disjunction uses exact cut elimination. For arithmetic,
\cref{thm:p07-proof-carrying-soundness} constructs proof components as
well as computational data. Ordinary realizability alone does not
supply the asserted HA proof.

\begin{hbconvention}{def:p07-ha-extraction-source}{Constructive extraction and its scope}
The HA disjunction and numerical-existence proofs use the explicitly
defined proof-carrying relation. Its soundness proof constructs every
component proof by finite HA inferences, including the proof of a
universal conclusion in the induction case. It does not infer HA
provability from ordinary realizability, nor require a normalization
theorem for an unstated arithmetic proof-term calculus.

Compare \cite[Vol.~I, Chapter~3, \S5.6]{TroelstraVanDalen1988}.
Here the extraction manipulates HA proof codes and programs. It is a
different proof, not a reconstruction of that source's normalization
calculus.

The propositional case follows from
\cref{thm:p07-lj-exact-cut-elimination}. Its last-rule step is
\cref{thm:p07-lj-disjunction-last-rule}. No complexity bound on
arithmetic witness extraction is asserted.
\end{hbconvention}

The chapter distinguishes Dialectica extraction, ordinary number
realizability, proof-carrying HA extraction, and modified realizability.
The formal extended-Church-thesis characterization is proved in the
explicit partial-term calculus, including the necessity of its
antecedent restriction. Proof mining is treated in the
structural-extensions chapter.
